\documentclass[11pt]{article}
\usepackage[margin=1in]{geometry}
\usepackage{amsmath}
\usepackage{amssymb}
\usepackage{hyperref}
\usepackage[numbers]{natbib}
\usepackage{booktabs}
\usepackage{multirow}

\hypersetup{
    colorlinks=true,
    linkcolor=blue,
    citecolor=blue,
    urlcolor=blue
}

\title{Daily Paper Review --- March 24, 2026\\
\large Beyond Single Tokens: Distilling Discrete Diffusion\\Models via Discrete MMD}
\author{Internbot --- Automated Research Review}
\date{March 24, 2026}

\begin{document}
\maketitle

\begin{abstract}
This report reviews \textbf{Discrete Moment Matching Distillation (D-MMD)}~\cite{dmmd}, which provides the first effective distillation framework for discrete diffusion models.
While continuous diffusion models have a rich ecosystem of distillation methods (consistency models, progressive distillation, moment matching), discrete diffusion models lack effective distillation---prior single-token approaches suffer from quality and diversity collapse.
D-MMD addresses this by matching \textbf{multi-token joint distributions} via kernel-based Maximum Mean Discrepancy (MMD), going ``beyond single tokens'' to capture inter-token correlations that are essential for coherent few-step generation.
The result: 4--16$\times$ speedup with distilled students that can \emph{surpass} their teachers.
This directly extends our CubiD~\cite{cubid} and Omni-Diffusion~\cite{omnidiff} reviews: D-MMD makes discrete diffusion generation practical by compressing hundreds of sampling steps into 1--4 steps.
Combined with the xLSTM distillation pipeline~\cite{xlstm-distill}, this establishes distillation as a unifying efficiency technique across both autoregressive and diffusion paradigms.
\end{abstract}

\section{Paper Summary}

\subsection{Overview}
\begin{itemize}
    \item \textbf{Title:} Beyond Single Tokens: Distilling Discrete Diffusion Models via Discrete MMD
    \item \textbf{Authors:} Emiel Hoogeboom, David Ruhe, Jonathan Heek, Thomas Mensink, Tim Salimans (Google DeepMind)
    \item \textbf{arXiv:} \href{https://arxiv.org/abs/2603.20155}{2603.20155}
    \item \textbf{Topics:} Discrete diffusion, distillation, moment matching, inference efficiency
\end{itemize}

\subsection{Core Problem}

A stark asymmetry exists in diffusion model distillation:
\begin{itemize}
    \item \textbf{Continuous diffusion} has consistency models, progressive distillation, DMD, moment matching, rectified flows---reducing sampling from hundreds of steps to 1--8 with minimal quality loss.
    \item \textbf{Discrete diffusion} lacks effective distillation. Prior attempts suffer from \textbf{quality collapse} and \textbf{diversity collapse}.
\end{itemize}

Why is discrete distillation hard?
\begin{enumerate}
    \item \textbf{Non-differentiable state space}: Discrete tokens $\{1, \ldots, K\}^D$ lack the smooth gradient structure of $\mathbb{R}^d$. Standard distribution matching breaks down.
    \item \textbf{Single-token independence}: Most discrete diffusion models (MDLM, SEDD) predict each token position independently. With many steps, correlations emerge iteratively. With few steps, \emph{independent predictions produce incoherent outputs}.
    \item \textbf{No continuous relaxation}: The forward process involves discrete corruption (masking, random replacement), making moment matching and distribution alignment fundamentally harder.
    \item \textbf{The step-reduction paradox}: Fewer steps forces each step to make larger ``jumps,'' requiring simultaneous prediction of multiple correlated tokens---something the standard independent parameterization cannot do.
\end{enumerate}

\subsection{Proposed Method}

\paragraph{Discrete Moment Matching Distillation.}
D-MMD trains a student discrete diffusion model to match the teacher's output distribution using a \textbf{kernel-based Maximum Mean Discrepancy} over discrete token sequences:
\[
\mathrm{MMD}^2(p, q) = \left\| \mathbb{E}_{x \sim p}[k(x, \cdot)] - \mathbb{E}_{y \sim q}[k(y, \cdot)] \right\|_{\mathcal{H}}^2
\]
where $k$ is a kernel function for categorical sequences and $\mathcal{H}$ is the associated Reproducing Kernel Hilbert Space.

The method adapts the continuous-domain moment matching framework of Salimans et al.\ (NeurIPS 2024) to discrete spaces.
In the continuous case, the core constraint is matching conditional expectations: $\mathbb{E}_g[\hat{x} \mid z_s] = \mathbb{E}_q[x \mid z_s]$.
In the discrete case, expectations over categorical variables lack the same geometric structure, so D-MMD replaces $L_2$ moment matching with kernel-based MMD.

\paragraph{Beyond Single Tokens.}
This is the central innovation.
\textbf{Single-token approaches} decompose the distillation loss into independent per-position terms:
\[
\mathcal{L}_{\mathrm{single}} = \sum_d D_{\mathrm{KL}}\!\left(p_{\mathrm{teacher}}^d(x^d \mid x_t) \;\|\; p_{\mathrm{student}}^d(x^d \mid x_t)\right)
\]
This cannot capture correlations between positions.
When generating ``the cat sat,'' single-token matching treats each word independently---it cannot learn that ``the'' and ``sat'' are correlated.

\textbf{D-MMD's multi-token approach} matches the \textbf{joint distribution} over multiple positions simultaneously:
\[
\mathcal{L}_{\mathrm{D\text{-}MMD}} = \mathrm{MMD}^2\!\left(p_{\mathrm{teacher}}(x^1, x^2, \ldots, x^D \mid x_t),\; p_{\mathrm{student}}(x^1, x^2, \ldots, x^D \mid x_t)\right)
\]
The kernel $k$ operates on \emph{full sequences} (or subsequences), so the loss captures how token choices at different positions interact.
Ablations confirm that matching over more tokens simultaneously leads to better generation quality.

\paragraph{Training and Student Architecture.}
The student maintains the \textbf{same backbone architecture} as the teacher (same Transformer, same parameters).
Distillation changes the \emph{training objective}, not the architecture---the student learns to produce high-quality outputs in fewer steps.
Training involves: (1) sample noisy states $x_t$ from the forward process, (2) generate clean predictions from both teacher and student, (3) compute D-MMD loss, (4) backpropagate through the student.

The MMD is estimated empirically using samples from both models:
\[
\widehat{\mathrm{MMD}}^2 = \frac{1}{n^2}\sum_{i,j} k(x_i, x_j) - \frac{2}{nm}\sum_{i,j} k(x_i, y_j) + \frac{1}{m^2}\sum_{i,j} k(y_i, y_j)
\]
where $x_i \sim$ teacher and $y_j \sim$ student.

\subsection{Key Results}

\paragraph{Efficiency Gains.}
D-MMD achieves \textbf{4--16$\times$ speedup} over teacher models:
\begin{itemize}
    \item Teachers typically require 50--1000 sampling steps for discrete diffusion.
    \item D-MMD students generate with \textbf{1--4 steps} while maintaining quality.
    \item With 4--8 steps, D-MMD students can \textbf{surpass their teachers}---distillation averages out per-timestep inconsistencies in the teacher, producing more coherent generation.
\end{itemize}

\paragraph{Multi-Token vs.\ Single-Token.}
The ablation on joint modeling sequence length is the paper's strongest empirical finding:
\begin{itemize}
    \item Single-token KL distillation: suffers quality/diversity collapse at few steps.
    \item Multi-token D-MMD: maintains both quality AND diversity.
    \item Increasing the number of jointly modeled tokens progressively improves generation quality, confirming that multi-token dependency capture is essential.
\end{itemize}

\paragraph{Comparison with Prior Discrete Distillation.}

\begin{center}
\begin{tabular}{lccc}
\toprule
\textbf{Method} & \textbf{Multi-Token?} & \textbf{Quality} & \textbf{Diversity} \\
\midrule
Exact Conditional Matching & No & Collapse & Collapse \\
Di4C (dim.\ correlations) & Partial & Moderate & Limited \\
CD4LM (consistency) & No & Good (w/ curriculum) & Good \\
DiMO (one-step) & No & Competitive & Limited \\
\textbf{D-MMD (this paper)} & \textbf{Yes} & \textbf{Surpasses teacher} & \textbf{Maintained} \\
\bottomrule
\end{tabular}
\end{center}

D-MMD is the only method that principled handles multi-token dependencies, and the only one where distilled students can surpass their teachers.

\paragraph{Generality.}
The paper includes Appendix D (``D-MMD for other discrete diffusion models''), demonstrating generalization beyond the primary setup to alternative discrete diffusion architectures (absorbing, uniform noise, etc.).

\subsection{Limitations}
\begin{itemize}
    \item \textbf{Kernel computation cost}: MMD scales quadratically with sample count and involves pairwise sequence comparisons, which is expensive for very long sequences.
    \item \textbf{Kernel hyperparameter sensitivity}: Bandwidth and kernel choice require careful tuning---too narrow captures only exact matches, too broad loses discriminative power.
    \item \textbf{Limited theoretical convergence guarantees}: The method works well empirically but formal guarantees are less established than in the continuous case.
    \item \textbf{Primarily categorical data}: Extension to other discrete structures may require additional adaptation.
\end{itemize}

\section{Follow-up Research Directions}

\subsection{D-MMD for CubiD: Efficient High-Dimensional Discrete Diffusion}

\paragraph{Gap.}
CubiD~\cite{cubid} achieves state-of-the-art discrete generation on 768-dim tokens but requires 256--512 iterative demasking steps.
D-MMD can distill discrete diffusion models to 1--4 steps.
However, CubiD's per-element masking creates a much larger discrete space ($L^{768}$ possibilities per position) than standard discrete diffusion ($K$ vocabulary entries per position), making the kernel computation and sample complexity challenges more severe.

\paragraph{Approach.}
Design a \textbf{factored kernel} for CubiD's high-dimensional discrete tokens that decomposes the full $768$-dimensional comparison into groups of dimensions:
\[
k(x, y) = \prod_{g=1}^{G} k_g(x^{(g)}, y^{(g)})
\]
where each group $g$ covers $768/G$ dimensions.
This reduces the per-comparison cost while preserving cross-dimensional correlations within groups.
The ``beyond single tokens'' aspect operates at two levels: across spatial positions (as in standard D-MMD) AND across dimension groups within each position (unique to CubiD).
Combined with IndexCache's cross-layer reuse~\cite{indexcache} for the Transformer backbone, this could reduce CubiD's 512-step generation to an effective $\sim$32--64 steps.

\paragraph{Feasibility.}
Product kernels are well-studied in the kernel methods literature and maintain MMD's theoretical guarantees.
CubiD's dimension-wise quantization to $L=8$ levels means each dimension group has a manageable vocabulary size.
The main risk is that factored kernels may miss long-range dimensional correlations, but CubiD's ablations suggest most information is captured by nearby dimension groups.
Initial experiments could distill CubiD-L (946M) on ImageNet-256 and measure gFID at 4/8/16 steps.

\subsection{D-MMD for Diffusion Language Models (Omni-Diffusion)}

\paragraph{Gap.}
Omni-Diffusion~\cite{omnidiff} uses masked discrete diffusion for text generation with per-token masking.
Its multi-step sampling is a bottleneck for real-time applications.
D-MMD is directly applicable: MDLM-style text diffusion models predict tokens independently per position (the exact limitation D-MMD addresses), and the multi-token MMD kernel can capture linguistic coherence (grammar, semantics) that single-token KL cannot.

\paragraph{Approach.}
Apply D-MMD to distill Omni-Diffusion's text generation component:
(1) Use the full Omni-Diffusion model as teacher (50+ steps);
(2) Train a D-MMD student with 4--8 steps using a sequence-level kernel that captures n-gram coherence;
(3) For the multimodal setting, design a \textbf{cross-modal kernel} that jointly matches text and image token distributions, ensuring that the distilled model preserves text-image alignment.
The multi-token aspect is especially important for text: a 4-step text diffusion student must predict entire coherent phrases per step, not individual words.

\paragraph{Feasibility.}
D-MMD's Appendix D already demonstrates generalization to alternative discrete diffusion architectures.
The kernel design for text is natural: n-gram overlap kernels, embedding similarity kernels, or learned kernels based on language model features.
Omni-Diffusion's text benchmarks provide ready-made evaluation.
A 4--8$\times$ speedup on text generation would make discrete diffusion competitive with autoregressive decoding for short-to-medium outputs.

\subsection{Unifying Autoregressive and Diffusion Distillation}

\paragraph{Gap.}
Our review series has covered two distillation paradigms:
(1) \textbf{Autoregressive distillation}: xLSTM distillation~\cite{xlstm-distill} transfers Transformer knowledge to sub-quadratic students via sparse KL (top-256 logits);
(2) \textbf{Diffusion distillation}: D-MMD transfers multi-step discrete diffusion knowledge to few-step students via kernel MMD.
Both use KL-family objectives but in different settings.
No work has unified these: can a single distillation framework transfer knowledge from an autoregressive teacher to a discrete diffusion student (or vice versa)?

\paragraph{Approach.}
Design a \textbf{cross-paradigm distillation} framework:
Use an autoregressive LLM as teacher (generating token-by-token) and a discrete diffusion model as student (generating via iterative demasking).
The D-MMD kernel compares the teacher's sequential output distribution with the student's parallel output distribution at the \emph{sequence level}---this is valid because both produce sequences of discrete tokens, regardless of the generation order.
The sparse top-$k$ logit technique from xLSTM distillation could reduce the kernel computation cost: instead of full-vocabulary kernels, operate only on the top-256 tokens per position.
This would enable leveraging the massive autoregressive LLM ecosystem as teachers for discrete diffusion students that offer parallel generation.

\paragraph{Feasibility.}
The sequence-level MMD kernel is paradigm-agnostic by design---it only requires samples from both models, not access to their internal generation process.
The main challenge is distribution mismatch: autoregressive models generate left-to-right with perfect local coherence, while diffusion models generate all positions simultaneously with global consistency.
A curriculum approach could help: start by matching short subsequences (where both paradigms are comparable) and gradually increase to full sequences.
Evaluation on standard text generation benchmarks (perplexity, MAUVE) would measure whether the diffusion student captures the autoregressive teacher's quality.

\section{Conclusion}

D-MMD resolves the fundamental limitation of discrete diffusion distillation: single-token independence assumptions that cause quality and diversity collapse at few sampling steps.
By matching multi-token joint distributions via kernel-based MMD, D-MMD produces students that achieve 4--16$\times$ speedup while maintaining---or even surpassing---teacher quality.
This is the missing piece for practical discrete diffusion: CubiD showed that discrete diffusion can scale to high-dimensional tokens; Omni-Diffusion showed it can handle multimodal generation; D-MMD makes both fast enough for deployment.

The paper also establishes \textbf{distillation as a unifying efficiency paradigm} across our review series:
xLSTM distillation compresses autoregressive Transformers into sub-quadratic architectures;
D-MMD compresses multi-step discrete diffusion into few-step generators;
both achieve near-perfect teacher recovery with 4--16$\times$ speedup.
The natural synthesis---cross-paradigm distillation from autoregressive teachers to discrete diffusion students---could combine the quality of massive LLMs with the parallel generation of diffusion models.

\begin{thebibliography}{9}
\bibitem{dmmd} Hoogeboom, E., Ruhe, D., Heek, J., Mensink, T., \& Salimans, T. (2026). Beyond Single Tokens: Distilling Discrete Diffusion Models via Discrete MMD. \textit{arXiv preprint arXiv:2603.20155}.

\bibitem{cubid} Wang, Y., et al. (2026). Cubic Discrete Diffusion: Discrete Visual Generation on High-Dimensional Representation Tokens. \textit{CVPR 2026}. arXiv:2603.19232.

\bibitem{omnidiff} (2026). Omni-Diffusion: Unified Multimodal Understanding and Generation with Masked Discrete Diffusion. \textit{arXiv preprint arXiv:2603.06577}.

\bibitem{xlstm-distill} Hauzenberger, L., et al. (2026). Effective Distillation to Hybrid xLSTM Architectures. \textit{arXiv preprint arXiv:2603.15590}.

\bibitem{indexcache} Bai, Y., et al. (2026). IndexCache: Accelerating Sparse Attention via Cross-Layer Index Reuse. \textit{arXiv preprint arXiv:2603.12201}.
\end{thebibliography}

\end{document}
